\subsection{Euler's rotation theorem: every rotation has an axis}\label{sec:rot:euler-s-rotation}
\begin{deriv}{$R\in\SO$ always fixes a line}
We need a vector $n$ with $Rn=n$, i.e.\ $\det(R-I)=0$. Using $RR\T=I$ and $\det R=\det R\T=1$:
\[
\det(R-I)=\det(R-RR\T)=\det R\,\det(I-R\T)=\det\big((I-R)\T\big)=\det(I-R)=(-1)^3\det(R-I).
\]
So $\det(R-I)=-\det(R-I)$, hence it is $0$: eigenvalue $1$ exists, with a unit eigenvector $\nh$. The plane perpendicular to $\nh$ is mapped to itself (dot products are preserved), and there $R$ acts as a $2\times2$ orthogonal matrix of determinant $1$, which is a planar rotation by some $\theta$. Therefore every 3-D rotation is ``turn by $\theta$ about one fixed axis $\nh$''. (The odd dimension is what makes this work. In 4-D a rotation can turn two planes independently with no fixed axis at all.)
\end{deriv}
